% Source: physical pp. 114--118 (printed pp. 91--95). Coverage: Section 17.3 from its heading through the paragraph immediately before Section 17.4; Eqs. (17.3.1)--(17.3.14), all unnumbered displays, citations [3,4], and two footnotes. Figures/tables/dividers: none. Corrected erratum: Eq. (17.3.7) begins at order \hbar, consistently with the loop expansion and Eq. (17.3.8).
\subsection[Renormalization: General Gauge Theories]{Renormalization: General Gauge Theories\footnotemark}
\footnotetext{This section lies somewhat out of the book's main line of development, and may be omitted in a first reading.}

The proof of the renormalizability of non-Abelian gauge theories in the
previous section relied on a `brute force' analysis of the possible terms
in the action of dimensionality four or less. But as we saw in Chapter 12,
this limitation on the dimensionality of terms in the action can be at best
a good approximation. The successful renormalizable quantum field theories
that are used to describe the strong, weak, and electromagnetic interactions
are almost certainly effective field theories, accompanied with terms of
dimensionality $d>4$; these terms are normally not observed because they are
suppressed by $4-d$ powers of some very large mass, perhaps of order
$10^{16}$--$10^{18}$ GeV. Gravitation too can be described by an effective
field theory, in which the Lagrangian density contains not only the
Einstein--Hilbert term $-\sqrt{g}R/16\pi G$, but also all scalars constructed
from four or more derivatives of the gravitational field. We need to show
that gauge theories of this sort, which are not renormalizable in the
power-counting sense, are nonetheless renormalizable in the modern sense
that the ultraviolet divergences are governed by the gauge symmetries in
such a way that there is a counterterm available to cancel every
infinity~\hyperref[ch17-ref-3]{[3]}.

For this purpose, let us return to the action $S[\chi,\chi^{\ddagger}]$
introduced in Section 15.9, taken as a function of independent fields
$\chi^n$ (including gauge and matter fields $\phi^r$ and ghost fields
$\omega^A$ as well as the non-minimal fields $\omega^{*A}$ and $h^A$),
together with their antifields $\chi_n^{\ddagger}$. In theories like quantum
gravity or Yang--Mills theories, that are based on a closed gauge algebra
with structure constants $f^C{}_{AB}$, this action is constrained to be of
the form
\begin{equation}
S=I[\phi]+\omega^A f_A^r[\phi]\phi_r^{\ddagger}
+\frac12\omega^A\omega^B f^C{}_{AB}[\phi]\omega_C^{\ddagger}
-h^A\omega_A^{*\ddagger},
\label{eq:17.3.1}
\end{equation}
where $I[\phi]$ is invariant under the infinitesimal gauge transformations
$\phi^r\longrightarrow\phi^r+\epsilon^A f_A^r[\phi]$. (As in Section 15.9,
the indices $r$, $A$, etc., include a spacetime coordinate, over which we
integrate in sums over these indices.)

We shall not limit ourselves here to actions of this form, but we shall
suppose that the local symmetries of the theory are imposed by requiring
that the action must obey some `structural constraints' on its antifield
dependence, of which Eq.~\eqref{eq:17.3.1} provides just one example. The
structural constraints are assumed to be linear, in the sense that if they
are satisfied for $S+\mathcal S_1$ and for $S+\mathcal S_2$, then for any
constants $\alpha_1$, $\alpha_2$ they are satisfied for
$S+\alpha_1\mathcal S_1+\alpha_2\mathcal S_2$. We also impose the quantum
master equation \eqref{eq:15.9.35}
\begin{equation}
(S,S)-2\ii\hbar\Delta S=0,
\label{eq:17.3.2}
\end{equation}
with $\Delta S$ defined by Eq.~\eqref{eq:15.9.34}, and the factor $\hbar$ now made
explicit as a `loop-counting' parameter, in the sense described in the
footnote of the previous section.

The action is taken as a power series in $\hbar$
\begin{equation}
S=S_R+\hbar S_1+\hbar^2S_2+\cdots,
\label{eq:17.3.3}
\end{equation}
where $S_R$ is an action of the same general form as $S$, but with all
coupling parameters replaced with finite renormalized values, and the $S_N$
are a set of infinite counterterms. The action $S$ is supposed to satisfy
the quantum master equation \eqref{eq:17.3.2} for all $\hbar$, so $S_R$
satisfies the classical master equation
\begin{equation}
(S_R,S_R)=0,
\label{eq:17.3.4}
\end{equation}
while the counterterms satisfy
\begin{equation}
(S_R,S_N)=-\frac12\sum_{M=1}^{N-1}(S_M,S_{N-M})+\ii\Delta S_{N-1}.
\label{eq:17.3.5}
\end{equation}

The counterterms $S_N$ are not by themselves sufficient to cancel the
ultraviolet divergences in loop graphs. As a generalization of the
conventional renormalization of fields, we also have to introduce a set of
renormalized fields and antifields, defined in terms of the original fields
and antifields by an arbitrary anticanonical transformation. An infinitesimal
anticanonical transformation may be defined in terms of an infinitesimal
generating functional $\delta F$ by Eq.~\eqref{eq:15.9.26}, so a sequence of
anticanonical transformations
$G(t)\longrightarrow G(t+\delta t)=G(t)+(F(t)\delta t,G(t))$ (where $G(t)$
is any functional of fields and antifields, and $F(t)$ is the generating
functional) leads to a finite canonical transformation
$G\longrightarrow\widetilde G\equiv G(1)$, with
\begin{equation}
\frac{\dd}{\dd t}G(t)=(F(t),G(t)),
\qquad G(0)=G.
\label{eq:17.3.6}
\end{equation}
If $F(t)$ is given by a power series
\begin{equation}
F(t)=\hbar tF_1+\hbar^2t^2F_2+\cdots,
\label{eq:17.3.7}
\end{equation}
then Eqs.~\eqref{eq:17.3.3}, \eqref{eq:17.3.6}, and \eqref{eq:17.3.7}
yield a transformed action
\begin{align}
\widetilde S={}&S_R+\hbar\bigl[S_1+(F_1,S_R)\bigr] \notag\\
&+\hbar^2\bigl[S_2+(F_1,S_1)+(F_2,S_R)
+\tfrac12(F_1,(F_1,S_R))\bigr]+\cdots.
\label{eq:17.3.8}
\end{align}
The question is whether we can use what freedom we have to choose the $F_N$
and $S_N$ so as to cancel all infinities arising from loop graphs.

As already noted, the first term $S_R$ in Eq.~\eqref{eq:17.3.3} is
automatically finite. Suppose that by cancellation of infinities with $S_M$
and $F_M$ for $M<N$ it has been possible to eliminate all infinities in the
terms $\Gamma_M$ of order $\hbar^M$ in the quantum effective action with
$M<N$. As we saw in the previous section, the Zinn--Justin equation (derived
here by setting $\chi_n^{\ddagger}=K_n+\delta\Psi/\delta\chi^n$) tells us in
this case that the infinite part $\Gamma_{N,\infty}$ of the term in the
quantum effective action of order $\hbar^N$ satisfies the condition
\begin{equation}
(S_R,\Gamma_{N,\infty})=0.
\label{eq:17.3.9}
\end{equation}
The field and antifield variables $\chi^n$ and $\chi_n^{\ddagger}$ are related
to the variables $\chi^n$ and $K_n$ by an anticanonical transformation, which
preserves all antibrackets, so the antibracket in Eq.~\eqref{eq:17.3.9} may be
calculated in terms of $\chi^n$ and $\chi_n^{\ddagger}$ instead of $\chi^n$
and $K_n$.

The condition \eqref{eq:17.3.5} satisfied by the counterterm $S_N$ is not the
same as the condition \eqref{eq:17.3.9} satisfied by $\Gamma_{N,\infty}$.
However, given any $S_N^0$ that satisfies Eq.~\eqref{eq:17.3.5}, we can find
a class of other solutions
\begin{equation}
S_N=S_N^0+S_N',
\label{eq:17.3.10}
\end{equation}
where $S_N'$ is arbitrary, except for the condition that $S_R+S_N'$, like
$S_R+S_N^0$, satisfies the same symmetry conditions as $S_R$, and that
\begin{equation}
(S_R,S_N')=0
\label{eq:17.3.11}
\end{equation}
so as not to invalidate Eq.~\eqref{eq:17.3.5}. The infinite part of the
$N$th-order term in the quantum effective action may therefore be written
\begin{equation}
\Gamma_{N,\infty}=S_{N,\infty}'+(F_{N,\infty},S_R)+X_{N,\infty},
\label{eq:17.3.12}
\end{equation}
where $X_N$ consists of terms from loop graphs, as well as from the term
$S_N^0$ and various terms in $\Gamma$ that involve $S_M$ and $F_M$ for
$M<N$. For instance, for $N=2$ Eq.~\eqref{eq:17.3.8} gives
\begin{align*}
X_2={}&S_2^0+2(F_1,S_1)+(F_1,(F_1,S_R))
+\text{two-loop terms involving only }S_R\\
&+\text{one-loop terms involving }S_R,\ S_1\text{ and }F_1.
\end{align*}
For our purposes the only thing we need to know about $X_N$ is that it does
not involve $S_N'$ or $F_N$, and that it is invariant under any linearly
realized global symmetries of $S_R$.

Now, because $(S_R,S_R)=0$, the operation $F\longmapsto(S_R,F)$ is
nilpotent; for all $F$,
\begin{equation}
(S_R,(S_R,F))=0.
\label{eq:17.3.13}
\end{equation}
Hence it follows from Eqs.~\eqref{eq:17.3.9} and
\eqref{eq:17.3.11}--\eqref{eq:17.3.13} that
\begin{equation}
(S_R,X_{N,\infty})=0.
\label{eq:17.3.14}
\end{equation}
Any term in $X_{N,\infty}$ of the form $(S_R,Y)$ may be cancelled in
Eq.~\eqref{eq:17.3.12} by choosing $F_{N,\infty}$ equal to $Y$. Thus the
space of possible remaining infinite terms in $X_{N,\infty}$ that need to
be cancelled by the counterterm $S_{N,\infty}'$ consist of those functionals
$X$ that satisfy $(S_R,X)=0$, counting as equivalent functionals that differ
only by terms of the form $(S_R,Y)$. In other words, the infinities in
$\Gamma_N$ that need to be cancelled by the counterterms $S_N'$ belong to
the \emph{cohomology} of the mapping $X\longmapsto(S_R,X)$.

The possible form of the counterterm $S_N'$ is limited by the requirement
that $S_R+S_N'$ must satisfy whatever structural constraints are imposed on
the action $S$. Thus we can complete the proof of renormalizability if we can
show that the cohomology of the mapping $X\longmapsto(S_R,X)$ consists only
of functionals that satisfy this structural constraint.

In the case of quantum gravity coupled to the Yang--Mills fields of a
semisimple gauge symmetry, the symmetries of the theory are implemented by
the structural constraint \eqref{eq:17.3.1}. In this case there is a
theorem~\hyperref[ch17-ref-4]{[4]} that states that the cohomology of the
mapping $X\longmapsto(S_R,X)$ (on the space of local functionals, rather
than spacetime-dependent functions, of ghost number zero) consists
\footnote{Strictly speaking, this is valid if one requires that the constant coefficients in $S$ do not take special values for which $S$ would be invariant under a larger group of local symmetries. For instance, this excludes the case $S=0$.}
of functionals $A[\phi]$ that are invariant under the gauge transformation
$\phi^r\longrightarrow\phi^r+\epsilon^A f_A^r[\phi]$ with structure constants
$f^C{}_{AB}$. Any $N$th-order infinity of this sort may be cancelled with a
counterterm $S_N'$ of the same form, so although these theories are not
renormalizable in the conventional power-counting sense, they are
renormalizable in the sense that all infinities can be eliminated by a
choice of parameters in the original bare action $I[\phi]$ and by a suitable
renormalization of fields and antifields.

In other theories the cohomology of the map $X\longmapsto(S_R,X)$ contains
additional terms. This does not necessarily require a weakening of the
structural constraints, because the additional terms in the cohomology may
not correspond to actual ultraviolet divergences. For instance, in gauge
theories with $U(1)$ factors the cohomology contains
terms~\hyperref[ch17-ref-4]{[4]} corresponding to a redefinition of the
action of the $U(1)$ gauge symmetry on the various fields of the theory,
which if infinite would require us to weaken the structural constraints by
leaving the normalization of the transformation functions $f_A^r[\phi]$ in
Eq.~\eqref{eq:17.3.1} arbitrary. But this infinity is forbidden by the same
soft-photon theorems that tell us that ratios of $U(1)$ couplings to various
fields (like the ratios of the various lepton charges in quantum
electrodynamics) are unaffected by radiative corrections. (See Section
10.4.) Where the extra terms in the cohomology do contain ultraviolet
divergences, it is necessary to weaken the structural constraint imposed on
$S$ in order that there should be a counterterm for every possible
ultraviolet divergence. It is not known whether this will always be
possible; if not, some theories may have to be rejected because of their
unremovable ultraviolet divergences.
